% --- 题目 ---
\problem[P222-1 (22-1)]{设随机变量 $\displaystyle X_1, X_2, \dots, X_n$ 独立同分布. 且 $\displaystyle X_1$ 的 $\displaystyle 4$ 阶矩存在. 记 $\displaystyle \mu_k = E(X_1^k) \ (k=1,2,3,4)$, 则根据切比雪夫不等式, 对任意 $\displaystyle \varepsilon > 0$, 有
\[ P\left\{\left|\frac{1}{n}\sum_{i=1}^n X_i^2 - \mu_2\right| \ge \varepsilon\right\} \le (\quad) \]
}
\begin{tasks}(4)
  \task $\displaystyle \frac{\mu_4 - \mu_2^2}{n\varepsilon^2}$.
  \task $\displaystyle \frac{\mu_4 - \mu_2^2}{\sqrt{n}\varepsilon^2}$.
  \task $\displaystyle \frac{\mu_2 - \mu_1^2}{n\varepsilon^2}$.
  \task $\displaystyle \frac{\mu_2 - \mu_1^2}{\sqrt{n}\varepsilon^2}$.
\end{tasks}
\ansat{368；【十年真题】 - 考点：大数定律、中心极限定理与切比雪夫不等式 - 1}
\vspace{4em}

% --- 题目 ---
\problem[P223-例]{设随机变量 $\displaystyle X$ 服从二项分布 $\displaystyle B(100, 0.2)$, 则利用中心极限定理可得 $\displaystyle P\{X \ge 12\}$ 的近似值为 ( \quad )}
\begin{tasks}(4)
  \task $\displaystyle \Phi(2)$.
  \task $\displaystyle 1-\Phi(2)$.
  \task $\displaystyle \Phi(0.5)$.
  \task $\displaystyle 1-\Phi(0.5)$.
\end{tasks}
\ansat{223；【方法探究】 - 考点：大数定律、中心极限定理与切比雪夫不等式 - 例}

% --- 题目 ---
\problem[P223-变式1 (03-3)]{设总体 $\displaystyle X$ 服从参数为 2 的指数分布, $\displaystyle X_1, X_2, \dots, X_n$ 为来自总体 $\displaystyle X$ 的简单随机样本, 则当 $\displaystyle n \to \infty$ 时, $\displaystyle Y_n = \frac{1}{n} \sum_{i=1}^n X_i^2$ 依概率收敛于 \underline{\hspace{4em}}.}
\ansat{368；【方法探究】 - 考点：大数定律、中心极限定理与切比雪夫不等式 - 变式1}
